\documentclass[11pt,fontset=fandol]{ctexart}
\usepackage[a4paper,margin=1in]{geometry}
\usepackage{amsmath,amsfonts,amssymb}
\usepackage{graphicx}
\usepackage{caption}
\usepackage[hidelinks]{hyperref}
\usepackage{microtype}

\setlength{\textfloatsep}{8pt plus 2pt minus 2pt}
\setlength{\floatsep}{6pt plus 2pt minus 2pt}
\setlength{\intextsep}{6pt plus 2pt minus 2pt}
\setlength{\abovecaptionskip}{4pt}
\setlength{\belowcaptionskip}{0pt}

\begin{document}

\noindent 本节中,$\kappa$ 表示膜通透性,$T$ 表示命中 target 的时间,
$\tau_{\mathrm{out}}$ 表示首次离开 corridor band 的时间,$\tau_{\mathrm{mem}}$ 表示首次
真正穿膜的时间。我们在 \texttt{peak1}、\texttt{valley} 和 \texttt{peak2} 三个窗口上
比较两类互补的量:一类是按 $T$ 归一化后的比值或百分比,另一类是以期望步数表示的绝对预算。
更具体地说,对任意事件时间 $\tau$,relative timing ratio 指
\[
  \frac{E[\tau \mid \tau < T,\ T \in W]}{E[T \mid T \in W]},
\]
也就是事件平均发生时刻占窗口 $W$ 内平均命中时间的比例;outside-time share 指命中时间
落在相应窗口内的轨道中,停留在 corridor band 外部的时间占 $T$ 的比例;outside budget
指命中之前停留在 corridor band 外部的期望步数;post-crossing budget 则指首次穿膜之后到
命中之前剩余的期望步数。本文中的 outside 指 corridor band 外部的并集,即 left-side
pocket 与 merged outer/right-side bulk。

\begin{figure}[!htbp]
  \centering
  \caption{上图给出对称 one-target baseline 的空间分区。下图把 $\kappa=0$ 与
    $\kappa=0.0040$ 的 $f(t)$ 画在同一时间轴上,并对曲线上标出的 \texttt{peak1}、
    \texttt{valley} 与 \texttt{peak2} 三个窗口嵌入堆叠预算柱;三段分别表示 left side、
    corridor 与 merged outer/right side 的归一化时间份额。曲线上的 $t_1$、$t_v$ 与
    $t_2$ 标记给出真实窗口中心,堆叠柱仅为保证可读性而做横向错位。}
  \label{fig:fig1-geometry-curve-budget-cn}
  \includegraphics[width=0.95\linewidth]{figures/fig1_geometry_curve_budget}
\end{figure}

图~\ref{fig:fig1-geometry-curve-budget-cn} 给出了 valley 与第二峰窗口比较所需的几何
基线。无论在 $\kappa=0$ 还是在 $\kappa=0.0040$ 下,以 \texttt{peak2} 为中心的命中时间
窗口分配到 corridor band 外部的预命中预算都高于 valley:outside share 在 control 中由
约 $48\%$ 升到 $56\%$,在 $\kappa=0.0040$ 时由约 $47\%$ 升到 $56\%$。堆叠柱进一步显示,
这个差异主要来自 left-side share 的上升:第二峰窗口中 left-side share 约为 $30\%$,而
merged outer/right-side share 反而低于 valley。这一分解表明,第二峰对应的延迟到达主要
来自命中前 corridor 外停留的增加,且主要由 left-side pocket 承担。

\clearpage

\begin{figure}[!htbp]
  \centering
  \caption{$\tau_{\mathrm{out}}$ 的 relative timing ratio,其中
    $E[\tau_{\mathrm{out}} \mid \tau_{\mathrm{out}} < T,\ T \in W]/E[T \mid T \in W]$ 为
    主量,也就是事件平均发生时刻占窗口 $W$ 内平均命中时间的比例。柱顶黑字给出
    relative timing ratio,柱底同色百分比给出 outside-time share,右上角注记给出对应的
    outside budget,也就是命中之前停留在 corridor band 外部的期望步数。三组数据分别
    对应 $\kappa=0$、$\kappa=0.0040$ 与边界点 $\kappa=0.0152$。}
  \label{fig:fig2-tau-out-budget-cn}
  \includegraphics[width=0.85\linewidth]{figures/fig2_tau_out_budget}
\end{figure}

图~\ref{fig:fig2-tau-out-budget-cn} 表明,\texttt{peak1} 在三种 $\kappa$ 下都保持
direct-corridor 型特征:ratio 约为 $0.76\sim0.79$,outside share 约为 $9\%\sim10\%$。
相比之下,\texttt{valley} 与 \texttt{peak2} 的 exit ratio 都较小,仅凭归一化 exit 时刻
并不能区分两者;差别主要体现在 outside budget 上。在 $\kappa=0$ 时,outside budget 由
约 $718$ 步增至约 $1311$ 步;在 $\kappa=0.0040$ 时由约 $705$ 步增至约 $1288$ 步。这说明
第二峰并非简单的几何 timing effect,而与命中前累积的 outside excursion 有关;结合图~1,
该对比在当前 case 中主要由 left-side pocket 承担。

\clearpage

\begin{figure}[!htbp]
  \centering
  \caption{$\tau_{\mathrm{mem}}$ 的 relative timing ratio,也就是首次穿膜平均发生时刻
    占窗口 $W$ 内平均命中时间的比例。柱顶黑字给出 relative timing ratio,柱底同色
    百分比给出膜穿越概率 $P(\tau_{\mathrm{mem}} < T \mid T \in W)$,右上角注记给出
    post-crossing budget,也就是首次穿膜之后到命中之前剩余的期望步数。图中的 `n/a'
    表示在命中之前没有发生 crossing,因此该 ratio 不定义。三组数据分别对应
    $\kappa=0$、$\kappa=0.0040$ 与边界点 $\kappa=0.0152$。}
  \label{fig:fig3-tau-mem-budget-cn}
  \includegraphics[width=0.85\linewidth]{figures/fig3_tau_mem_budget}
\end{figure}

图~\ref{fig:fig3-tau-mem-budget-cn} 把膜通透性的修正单独抽离出来。在 $\kappa=0$ 时,
crossing probability 和 post-crossing budget 均为零。$\kappa>0$ 后,valley 与
\texttt{peak2} 的 crossing probability 仍较接近,例如在 $\kappa=0.0040$ 时约为
$18\%$ 对 $20\%$;更明显的差异出现在首次穿膜后的剩余演化时间上,约为 $213$ 步对
$336$ 步。到边界点 $\kappa=0.0152$ 时,post-crossing budget 进一步增至约 $604$ 步与
$785$ 步。这说明膜通透性主要增大首次穿膜之后到命中之前的期望步数,从而放大晚时间
shoulder/second-peak 的贡献;crossing probability 本身不足以解释第二峰。

\end{document}
